\section{What a semantic layer has to represent before it can help}
\label{sec:product}

The negative governed result is more useful as a specification than as a verdict.
Each item below names a capability whose absence is visible in the preserved
artifacts, and each is a general property of importing business knowledge into a
governed model rather than an observation about one benchmark.

\paragraph{Grain, identity, cardinality, and aggregation need to be first-class,
and their absence needs to be visible before deployment.}
Table~\ref{tab:hkb} is the evidence. A conservative compiler that refuses to
guess turned 17.7\% of definitions into executable objects and deferred 46.9\% of
them across an unresolved grain, and the three most common recorded reasons were
unknown cardinality, unspecified aggregation, and missing cross-grain identity.
Those are the four contracts an import path currently has no way to express. The
practical consequence is that a model author cannot tell, before deploying,
whether the model carries enough structure to govern anything. A dry run that
reports which definitions compiled, which were kept as context, which were
deferred, and why, converts a silent coverage problem into a decision. Without
it, the choice is between silently guessed semantics and a large volume of prose
whose discoverability by the agent is unknown.

\paragraph{Relationship coverage is the same problem one level up.}
The public schemas of these 16 databases contain 1,049 foreign keys that pass a
conservative cardinality contract, and the bounded modeled candidate could expose
91 relationships. A model author needs to see that gap, with accepted, deferred,
and unreachable relationships enumerated, before deciding whether the model has
enough structure for a governed query path to exist at all. In this study the gap
was not visible until the query-path measurement showed that no governed query
had ever declared a join.

\paragraph{Retrieval payloads belong in the agent contract.}
The first direct-SQL attempt in this study returned all 51 tables of a database
and then exhausted its budget before writing a line of SQL
(Table~\ref{tab:slice}). A semantically valid tool call ended the run. Schema and
semantic search should be query-directed, bounded, and observable, with per-call
context volume in telemetry; and typed public provenance identifiers should exist
so that a safety filter does not reject canonical foreign-key values as
suspected secrets, which is the second failure the same slice exposed.

\paragraph{Not answering has several causes, and the product currently reports
one.}
A content refusal, an explicit statement that the available schema is
insufficient, an exhausted model budget, and an infrastructure failure imply four
different product actions. The governed condition reports them as one terminal
state, which is why its refusal rate is unavailable in Table~\ref{tab:efficiency}
rather than estimated. Governed jobs should expose a typed refusal outcome, and
both the product and comparator telemetry should retain the raw state even when a
report also groups them into broader reliability categories.

\paragraph{The result contract over rewritten SQL is the one nobody wrote down.}
Thirty-one of the 34 governed non-answers in the development baseline shared an
\texttt{UNKNOWN} selected-field type, and the relationship experiment then
preserved 14 generated semantic queries that could not be captured because their
result types were unsupported. That is not a scoring inconvenience; it is a gap
in the interface. When the agent authors the query, neither side currently
specifies what the planner guarantees about the type of an output column that the
semantic model never declared. Production planning and execution need a stable,
total representation for unknown, boolean, temporal, and null values, and an
unsupported planner type should surface as a visible product outcome rather than
an adapter exception that obscures whether the governed query itself succeeded.
The same argument covers truncation: a preview is not an analytical result, and
presentation-control records do not belong in data rows.

\paragraph{Ordering.}
These are not independent. Result typing comes first, because until an unsupported
output type stops silently consuming a governed attempt, no downstream experiment
can be scored cleanly, as E02's 14 uncaptured queries demonstrate. Grain and
relationship authoring comes next, because it is what would give a planner
something to compile. Retrieval bounds, refusal typing, and deployment
diagnostics are independent of both and cheaper than either.

The registered C5 condition tested the second of these directly. It published the
joins and the full knowledge context that the mechanical baseline lacked, and the
query-path readout in Table~\ref{tab:c5mechanism} settles the question:
relationship coverage alone does not move a governed agent off the rewrite path,
because not one of the 134 parseable C5 queries left it. Accuracy nevertheless
doubled, so the full knowledge context is doing real work as context rather than
as compiled semantics, and the grain and aggregation contracts remain the
untested half. Its deployment
history already exercises two of the cheaper items: the readback divergence in
Section~\ref{sec:c5} is the canonical-export gap, and the branch that had to be
uploaded twice is the deployment-identity gap.
